% Post-Phase-18 extension, activated by explicit user approval on 2026-09-03.
\section{Extension: Resource Allocation Between Research and Development}
\label{sec:erd-extension}

\subsection{Separation from the commercialization baseline}
\label{subsec:erd-separation}

This section activates only the research-versus-development block previewed
in Section~\ref{subsec:p11-research-development}. It does not alter the
commercialization baseline. In particular, the institutional regime still
acts directly only through \(M\mapsto\tau_E(M)\), route choice remains
deterministic, and \(\Omega_i(M,p_m)\) continues to be the expected optimized
value per route-planning-stage project supplied by the baseline route block.
The baseline control \(x_i\) is not retroactively redefined. When this
extension is activated, it is replaced by two extension-only controls:
upstream research allocation \(x_i^R\) and development / project-advancement
allocation \(x_i^D\). Thus \(x_i^D\), rather than \(x_i^R\), is the analogue
of baseline project advancement.

The extension addresses a distinct question. It asks how an increase in
anticipated commercialization value reallocates a finite innovation resource
between research and development, and when that reallocation can make patent
applications fall even though development rises. It does not make patenting
a direct policy target or a direct component of the MAH route wedge.

\subsection{Technology, resource constraint, and objective}
\label{subsec:erd-problem}

Both allocations are measured in normalized innovation-resource units
\(\mathsf X\). Developer \(i\) has a resource ceiling
\(\bar X_i>0\). The expected measure of projects reaching the route-planning
stage is
\begin{equation}
  N_i(x_i^R,x_i^D)
  =
  \left(a_i+\gamma_i x_i^R\right)x_i^D,
  \qquad \gamma_i\geq0,
  \label{eq:erd-project-technology}
\end{equation}
where \(a_i>0\) remains predetermined and \(\gamma_i\) measures
research--development complementarity. The extension does not permit a
direct policy shift in either object.

For a fixed candidate CMO price and hence a fixed value
\(\Omega_i=\Omega_i(M,p_m)\), the developer solves
\begin{equation}
  \begin{aligned}
  \max_{x_i^R,x_i^D\geq0}\quad
  \mathcal V_i
  ={}&
  u_i x_i^R-\frac{\kappa_{Ri}}{2}(x_i^R)^2
  +\beta\Omega_i
       \left(a_i+\gamma_i x_i^R\right)x_i^D
  -\frac{\kappa_{Di}}{2}(x_i^D)^2,\\
  \text{subject to}\quad
  &x_i^R+x_i^D\leq\bar X_i,
  \end{aligned}
  \label{eq:erd-objective}
\end{equation}
where \(u_i>0\), \(\kappa_{Ri}>0\), and \(\kappa_{Di}>0\) are policy-invariant.
The first two terms summarize the private value and convex organizational
cost of upstream research. The fourth term is the convex organizational cost
of development. The resource ceiling is separate from those real costs: it
captures financing or managerial resources that cannot be expanded within
the decision cohort. Consequently, a positive multiplier on the ceiling,
not merely a low observed resource proxy, defines a constrained firm in the
model.

Let
\begin{equation}
  t_i=\beta\Omega_i,
  \qquad
  \Delta_i(t_i)
  =\kappa_{Ri}\kappa_{Di}-(t_i\gamma_i)^2.
  \label{eq:erd-concavity-index}
\end{equation}
The maintained extension condition is \(\Delta_i(t)>0\) for every value
\(t=\beta\Omega_i(M,p_m)\) reached in the comparison. This makes the Hessian
of \(\mathcal V_i\) negative definite. Together with the compact feasible
set, it yields a unique optimizer. This restriction is an extension-specific
curvature condition, not a baseline assumption.

\subsection{KKT system}
\label{subsec:erd-kkt}

Let \(\varpi_i\geq0\) be the multiplier on the common resource ceiling, and
let \(\zeta_i^R,\zeta_i^D\geq0\) be the nonnegativity multipliers. The complete
KKT system is
\begin{align}
  u_i-\kappa_{Ri}x_i^R+t_i\gamma_i x_i^D
  -\varpi_i+\zeta_i^R&=0,
  \label{eq:erd-kkt-research}\\
  t_i(a_i+\gamma_i x_i^R)-\kappa_{Di}x_i^D
  -\varpi_i+\zeta_i^D&=0,
  \label{eq:erd-kkt-development}\\
  \varpi_i(\bar X_i-x_i^R-x_i^D)&=0,
  \label{eq:erd-kkt-budget}\\
  \zeta_i^R x_i^R=\zeta_i^D x_i^D&=0,
  \label{eq:erd-kkt-nonnegative}
\end{align}
together with primal and dual feasibility. Strict concavity makes these
conditions sufficient as well as necessary.

The sign results below concern interior allocations. At a nonnegativity
corner, the corresponding strict sign becomes weak and is governed by the
one-sided KKT inequality. If a reform moves a firm across the slack/binding
boundary, the two regime formulas must be evaluated as a finite comparison;
one may not differentiate through the kink.

\subsection{Slack resources: development rises and research does not fall}
\label{subsec:erd-slack}

Suppose the unconstrained optimizer is interior and satisfies
\(x_i^R+x_i^D<\bar X_i\). Then \(\varpi_i=0\), and the unique solution is
\begin{align}
  x_i^{R,u}(t_i)
  &=
  \frac{\kappa_{Di}u_i+t_i^2\gamma_i a_i}
       {\Delta_i(t_i)},
  \label{eq:erd-slack-research}\\
  x_i^{D,u}(t_i)
  &=
  \frac{t_i(\gamma_i u_i+\kappa_{Ri}a_i)}
       {\Delta_i(t_i)}.
  \label{eq:erd-slack-development}
\end{align}

\begin{extensionproposition}[Slack-resource response]
\label{prop:erd-slack}
Hold the candidate CMO price fixed and vary the continuous argument
\(\Omega_i\). At an interior slack solution satisfying
\(\Delta_i(t_i)>0\),
\begin{align}
  \frac{\partial x_i^{R,u}}{\partial\Omega_i}
  &=
  \frac{2\beta t_i\gamma_i\kappa_{Di}
        (a_i\kappa_{Ri}+\gamma_i u_i)}
       {\Delta_i(t_i)^2}
  \geq0,
  \label{eq:erd-slack-research-cs}\\
  \frac{\partial x_i^{D,u}}{\partial\Omega_i}
  &=
  \frac{\beta(\gamma_i u_i+\kappa_{Ri}a_i)
        [\kappa_{Ri}\kappa_{Di}+t_i^2\gamma_i^2]}
       {\Delta_i(t_i)^2}
  >0.
  \label{eq:erd-slack-development-cs}
\end{align}
The research response is strict when \(\gamma_i>0\), and is zero when
\(\gamma_i=0\).
\end{extensionproposition}

\begin{proof}
With \(\varpi_i=\zeta_i^R=\zeta_i^D=0\), equations
\eqref{eq:erd-kkt-research}--\eqref{eq:erd-kkt-development} are a two-by-two
linear system with coefficient determinant \(\Delta_i(t_i)>0\). Inversion
gives \eqref{eq:erd-slack-research}--\eqref{eq:erd-slack-development}.
Differentiating those expressions with respect to \(t_i\), then using
\(\partial t_i/\partial\Omega_i=\beta\), gives
\eqref{eq:erd-slack-research-cs}--\eqref{eq:erd-slack-development-cs}.
All factors have the stated signs. The result is a derivative with respect to
the continuous value argument \(\Omega_i\), not a derivative with respect to
binary policy \(M\).
\end{proof}

Thus an unconstrained firm does not need to finance more development by
reducing research. If research raises the productivity of development, both
allocations rise. If that complementarity is absent, development rises while
research remains locally unchanged.

\subsection{Binding resources: development can crowd out research}
\label{subsec:erd-binding}

Suppose instead that the unique optimizer is interior and the resource
constraint binds. This occurs when the unconstrained demand for the two
activities exceeds \(\bar X_i\), subject to the resulting constrained
allocation remaining interior. Using \(x_i^D=\bar X_i-x_i^R\) and equating
the two interior marginal conditions yields
\begin{align}
  x_i^{R,b}(t_i)
  &=
  \frac{u_i+\kappa_{Di}\bar X_i
        +t_i(\gamma_i\bar X_i-a_i)}
       {\kappa_{Ri}+\kappa_{Di}+2t_i\gamma_i},
  \label{eq:erd-binding-research}\\
  x_i^{D,b}(t_i)&=\bar X_i-x_i^{R,b}(t_i).
  \label{eq:erd-binding-development}
\end{align}

\begin{extensionproposition}[Binding-resource reallocation]
\label{prop:erd-binding}
At an interior binding solution,
\begin{align}
  \frac{\partial x_i^{R,b}}{\partial\Omega_i}
  &=
  \frac{\beta[\gamma_i(x_i^{D,b}-x_i^{R,b})-a_i]}
       {\kappa_{Ri}+\kappa_{Di}+2t_i\gamma_i},
  \label{eq:erd-binding-research-cs}\\
  \frac{\partial x_i^{D,b}}{\partial\Omega_i}
  &=-\frac{\partial x_i^{R,b}}{\partial\Omega_i}.
  \label{eq:erd-binding-development-cs}
\end{align}
Therefore, under the development-biased value-shift condition
\begin{equation}
  a_i>\gamma_i(x_i^{D,b}-x_i^{R,b}),
  \label{eq:erd-development-biased-condition}
\end{equation}
an increase in anticipated commercialization value strictly raises
development and strictly reduces research. If equality holds, there is no
local reallocation. If the inequality is reversed, the allocation signs are
reversed and the extension does not predict a patent decline.
\end{extensionproposition}

\begin{proof}
Differentiate the difference between the two interior KKT equations while
using \(dx_i^D=-dx_i^R\). This gives
\[
  [\kappa_{Ri}+\kappa_{Di}+2t_i\gamma_i],dx_i^R
  =
  [\gamma_i(x_i^D-x_i^R)-a_i],dt_i.
\]
Substituting \(dt_i=\beta\,d\Omega_i\) proves
\eqref{eq:erd-binding-research-cs}--\eqref{eq:erd-binding-development-cs}.
The denominator is strictly positive, so condition
\eqref{eq:erd-development-biased-condition} determines the signs.
\end{proof}

Although the two allocations move in opposite directions, the optimal
planning-stage project mass is weakly increasing in commercialization value.
At an interior binding allocation,
\begin{equation}
  \frac{\partial N_i^*}{\partial\Omega_i}
  =
  \frac{\beta[a_i-\gamma_i(x_i^{D,b}-x_i^{R,b})]^2}
       {\kappa_{Ri}+\kappa_{Di}+2t_i\gamma_i}
  \geq0.
  \label{eq:erd-binding-project-mass}
\end{equation}
This identity separates the development response from the research response:
project advancement can rise even when upstream research falls.

\subsection{Patent applications and the evidence boundary}
\label{subsec:erd-patents}

To create an explicit patent margin, let the expected flow of patent
applications be
\begin{equation}
  P_i^A=h_i(x_i^R),
  \qquad h_i'(x_i^R)>0.
  \label{eq:erd-patent-production}
\end{equation}
This is a measurement/production equation for the extension. It does not say
that every unit of research produces a patent, that every patent has equal
value, or that observed applications identify \(x_i^R\) without error.

\begin{extensioncorollary}[Patent sign]
\label{cor:erd-patent-sign}
Consider two policy regimes whose equilibrium route values satisfy
\(\Omega_i^1>\Omega_i^0\).

\begin{enumerate}
  \item If the firm is interior and resource-slack throughout the comparison,
  then \(P_i^{A,1}\geq P_i^{A,0}\), with a strict inequality when
  \(\gamma_i>0\).
  \item If the firm is interior and resource-constrained throughout the
  comparison and condition
  \eqref{eq:erd-development-biased-condition} holds along the comparison
  path, then
  \(x_i^{D,1}>x_i^{D,0}\) and \(P_i^{A,1}<P_i^{A,0}\).
\end{enumerate}

If the resource regime changes between the two policy states, if the net
value gain is zero, or if the development-biased condition fails, the patent
sign must be computed from the finite solutions and is not signed by this
corollary.
\end{extensioncorollary}

The evidence that granted or core patents do not significantly decline is
also compatible with, but not forced by, this extension. For example, write
\(P_i^A=P_i^H+P_i^L\). If the high-value component is locally inelastic to
the marginal research reallocation while \(dP_i^L/dx_i^R>0\), the decline is
concentrated in low-value applications and \(P_i^H\) need not change. This is
an additional measurement/composition condition. It is not a theorem that
all high-value patents are protected.

\subsection{CMO-price feedback and the binary MAH comparison}
\label{subsec:erd-cmo-feedback}

The policy remains binary. Its effect must therefore be evaluated by the
finite equilibrium value comparison, not by differentiating with respect to
\(M\). The extension changes study-related CMO demand because the planning
mass is now \(N_i^*\):
\begin{equation}
  D_{m,RD}^{\mathrm{MAH}}(p_m;M)
  =
  \int
  N_i^*(M,p_m)\chi_i^E(p_m;M)\,dH(a,k),
  \label{eq:erd-cmo-demand}
\end{equation}
where \(\chi_i^E\) is the baseline expected qualified capacity per
planning-stage project. Supplier behavior, background demand, and the single
market-clearing condition are unchanged.

The optimized \(N_i^*\) is nondecreasing in \(\Omega_i\), including across a
slack/binding switch. To see this without differentiating through the kink,
apply revealed preference to two values \(\Omega_i'<\Omega_i''\). Adding the
two optimality inequalities gives
\[
  \beta(\Omega_i''-\Omega_i')
  [N_i^*(\Omega_i'')-N_i^*(\Omega_i')]
  \geq0.
\]
Because the baseline envelope gives \(\partial\Omega_i/\partial p_m
=-\chi_i^E\leq0\), and \(\chi_i^E\) is also weakly decreasing in the
candidate price, equation \eqref{eq:erd-cmo-demand} remains weakly decreasing
in \(p_m\) under the baseline aggregation regularity. Strictly increasing
supply therefore continues to select at most one CMO price. The extension
adds no second market-clearing condition.

Let \(p_{m,RD}^h\) denote the resulting CMO equilibrium in regime
\(h\in\{0,1\}\), and define the net equilibrium value gain
\begin{equation}
  \Delta\Omega_{i,RD}^{\mathrm{eq}}
  =
  \Omega_i(1,p_{m,RD}^1)
  -\Omega_i(0,p_{m,RD}^0).
  \label{eq:erd-equilibrium-value-gain}
\end{equation}
The baseline choice-set argument implies
\(\Delta\Omega_{i,RD}^{\mathrm{eq}}\geq0\); scarcity can attenuate the gain
and can reduce it to zero. Propositions~\ref{prop:erd-slack} and
\ref{prop:erd-binding}, and Corollary~\ref{cor:erd-patent-sign}, map into an
MAH prediction only when this net gain is strict and their resource-regime
conditions hold.

\subsection{Interpretation for Gu (2024) and remaining non-predictions}
\label{subsec:erd-gu-mapping}

The extension rationalizes the joint pattern documented by \citet{Gu2024}
without changing the commercialization baseline. Firms with a positive
resource multiplier can respond to a strict MAH-induced commercialization
value gain by shifting a fixed resource toward development, raising
development-stage activity while reducing upstream research and patent
applications. Firms with slack resources raise development without crowding
out research; when \(\gamma_i=0\), their research and patent applications are
locally unchanged. This creates a sharp resource-heterogeneity prediction.
Observed balance-sheet resource groups remain proxies, however: they do not
directly reveal \(\varpi_i\).

The original/incremental class \(g\in\{O,\mathrm{Inc}\}\) remains an exogenous
classifier with common developer-level allocations \(x_i^R,x_i^D\). The
allocation block imposes no class ranking. Class-specific route and realized
outcomes continue to use \(\rho_gF_g\); a stronger incremental response
requires the incremental conditional distribution to place more mass on
projects receiving a strict route-value or observed-outcome gain. The
extension therefore accommodates, but does not assume, the class pattern in
Gu (2024). Likewise, transition from a pure researcher into development is
an extensive-margin entry/state-transition question and remains outside this
fixed-population extension.

The principal boundary cases are:

\begin{itemize}
  \item if \(\Delta\Omega_{i,RD}^{\mathrm{eq}}=0\), MAH causes no allocation
  or patent response through this channel;
  \item if \(\gamma_i=0\), slack-firm research is unchanged, whereas a
  binding firm reallocates toward development whenever \(a_i>0\);
  \item as \(\bar X_i\) becomes sufficiently large, the resource constraint
  becomes slack and the crowd-out result disappears;
  \item at a nonnegativity corner or a slack/binding regime switch, strict
  local formulas are replaced by KKT inequalities or direct finite
  comparison;
  \item if \(h_i'=0\) over the observed range, research can change without a
  measured patent-count response; and
  \item failure of the curvature or development-biased condition removes the
  corresponding sign claim rather than reversing it by assumption.
\end{itemize}

Accordingly, the model now has an explicit patent prediction in this
extension, but not an unconditional one: the sign depends on whether the
common resource constraint binds and on the direction of the marginal value
shift.
